% What the Cache Forgets / Sharpens -- ACL long paper (8 pages + limitations + appendix).
% Build: bash build_paper.sh (repo root) -> paper.pdf and course_paper.pdf.
% EVERY result number comes from ../../tables/*.tex (written by make_results.py); nothing is typed by hand.
% The pre-registered verdict was inconclusive (RESULTS.md); the text below states it directly.
\documentclass[11pt]{article}

% ---- where the generated tables/figures live (repo root when built in place; ./ in the Overleaf zip)
\makeatletter
\IfFileExists{tables/macros.tex}{\def\resultsdir{}}{\def\resultsdir{../../}}
\makeatother

% ---- ACL style if present (build_paper.sh fetches it); otherwise a close article fallback
\IfFileExists{acl.sty}{%
  \usepackage[final]{acl}%
}{%
  \usepackage[margin=2.5cm]{geometry}%
  \usepackage[round]{natbib}%
  \usepackage{times}%
}
\usepackage{times}
\usepackage{latexsym}
\usepackage[T1]{fontenc}
\usepackage[utf8]{inputenc}
\usepackage{microtype}
\usepackage{amsmath,amssymb}
\usepackage{booktabs}
\usepackage{graphicx}
\usepackage{makecell}
\usepackage{adjustbox}
\usepackage{xcolor}
\usepackage{url}

% course version: build_paper.sh compiles with \def\COURSE{1}
\newif\ifcourse
\ifdefined\COURSE \coursetrue \fi

\input{\resultsdir tables/macros}

\newcommand{\KL}{\mathrm{KL}}
\newcommand{\figorbox}[2]{\IfFileExists{\resultsdir figures/#1.pdf}{\includegraphics[width=#2]{\resultsdir figures/#1.pdf}}{\fbox{\parbox{0.9\columnwidth}{\centering\small figure #1: not available yet}}}}

\title{Does the Cache Forget? Top-$k$ Teacher Caching Sharpens Distilled Students but Loses Tail Identity\ifcourse\\\normalsize Course project report\fi}

\author{Sabuhi Aghayev \and Ilkin Mammadli}

\begin{document}
\maketitle

% =====================================================================================================
\begin{abstract}
Distilling a large language model on limited compute usually means caching only the teacher's top-$k$ next-token
probabilities. We ask what this routine choice does to what the student \emph{generates}. Distilling
SmolLM2-1.7B-Instruct into SmolLM2-135M on Dolly with seven teacher-signal formats, three training seeds and a
pre-registered matched-temperature test, we find that top-$k$ caching makes students sharper: at the default
sampling temperature a top-8 student is \SemDivDropPct\% less diverse, has higher reference overlap and teacher
compatibility, and has $\times$\ECERatio{} the next-token calibration error of full distillation, at statistically
equivalent greedy ROUGE-L, monotonically in $k$. This is not a fixed loss of diversity: at matched reference overlap
the pre-registered test does not support a residual deficit, and the point estimates favour the top-8 student. What
caching loses is the identity of the tail. Restoring only the amount of tail mass, with a ghost token or even spread
uniformly, recovers diversity and calibration but not within-tail KL, and only partly rare-token prediction;
random-sampling distillation recovers both. All claims are for one model family and one training set.
\end{abstract}

% =====================================================================================================
\section{Introduction}

Knowledge distillation (KD) trains a small \emph{student} to imitate a larger \emph{teacher}
\citep{hinton2015distilling}. For language models, token-level KD matches the student's next-token distribution to
the teacher's at every position of a training sequence. When compute is limited the teacher cannot run next to
the student at every step, so its outputs are computed once and \emph{cached}. A full distribution over a
49k-token vocabulary is too large to store for millions of positions, so practitioners keep only the teacher's
top-$k$ probabilities and renormalise them, with $k$ chosen for convenience \citep{offlinetopk2026,topkgradient2026}.

Prior work established that this truncation is biased: renormalised top-$k$ targets are not an unbiased estimate of
the teacher \citep{anshumann2025sparse}, and they give no gradient to the student's probability mass outside the
cached set, so that mass drifts \citep{topkgradient2026}. These studies measure loss, accuracy and calibration.
What the student \emph{generates} has not been measured. A student trained to put all its mass on the teacher's
top-$k$ tokens could become systematically sharper: less diverse in what it writes, less faithful to the teacher's
uncertainty, and worse on rare tokens that live in the teacher's tail. If so, top-$k$ caching is a hidden source
of the output homogenisation that recent work attributes to distillation \citep{quantdistill2025}.

We study this directly. We distil SmolLM2-1.7B-Instruct into SmolLM2-135M on databricks-dolly-15k with seven
teacher-signal formats (gold-token SFT, sequence-level KD, renormalised top-$k$ for $k\in\{8,32,128\}$, top-$k$
plus a tail bucket, random-sampling KD, and full-vocabulary KD), each with forward and reverse KL where defined,
identical data order, steps and learning rate, and three training seeds. Our contributions:
\begin{enumerate}
\item \textbf{A measurement of what top-$k$ caching changes in generation.} Top-8 caching reduces sampled
  diversity and worsens calibration despite similar greedy ROUGE-L, with a monotone dose--response in $k$
  (\S\ref{sec:main}); the effect survives length controls (\S\ref{sec:length}).
\item \textbf{A matched-temperature test, pre-registered.} Because top-8 samples are sharper, we select sampling
  temperatures on validation so that both students have equal reference overlap and compare diversity on test
  with a paired bootstrap and a decision rule fixed before the runs (\S\ref{sec:temp}). A residual diversity deficit
  is not supported: at matched reference overlap the top-8 student is, if anything, more diverse.
\item \textbf{Tail amount vs.\ tail identity.} Using the decomposition of the KL from \citet{dasgupta2026tad}, we
  show that restoring how much mass the student puts outside the cache (ghost token, or a uniform spread) does not
  restore which tail tokens receive it; RS-KD's approximation fidelity extends to the generation endpoints measured
  (\S\ref{sec:tail}).
\item \textbf{Open, reproducible code on one 8\,GB GPU}, with every run logged, including failures (reverse KL
  from the base checkpoint does not terminate; \S\ref{sec:rkl}).
\end{enumerate}

% =====================================================================================================
\section{Background and related work}

\paragraph{Token-level KD and the divergence.} Forward KL, $\KL(p\Vert q)$, is mass-covering; reverse KL,
$\KL(q\Vert p)$, is mode-seeking and was proposed for LLM distillation by MiniLLM \citep{gu2024minillm}; DistiLLM
uses a skewed KL \citep{ko2024distillm}. Several works report that reverse KL reduces output diversity and add
forward-KL terms or modify the gradient \citep{jin2026eopd,luong2026drkl}. These studies vary the divergence; we
vary the \emph{format of the cached teacher signal} and its interaction with the divergence.

\paragraph{Sparse and cached-logit KD.} \citet{anshumann2025sparse} show that caching top-$k$ probabilities biases
the teacher estimate and harms calibration, and propose Random Sampling KD (RS-KD): draw tokens from the teacher and
train on their empirical frequencies, unbiased in expectation. They also define a ``ghost token'' that collects
the teacher's leftover mass as one extra class; the same $(k{+}1)$-way construction appears in TAD
\citep{dasgupta2026tad} and, for on-policy reverse KL, in TA-OPD \citep{huang2026taopd}. \citet{topkgradient2026}
prove that renormalised top-$k$ KD gives zero gradient along tail-controlling directions and show tail-mass drift.
None of these works measures output diversity, the tail coverage of generations, or long-tail behaviour; this is
the gap we address. We use the ghost token and RS-KD as baselines, not as contributions.

\paragraph{Homogenisation and long-tail knowledge.} Distillation homogenises models \citep{quantdistill2025}, and
long-tail facts are learned poorly \citep{kandpal2023longtail,mallen2023popqa}. We probe token-level tail behaviour
and, exploratorily, factual recall by entity popularity.

% =====================================================================================================
\section{Method}\label{sec:method}

Let $x$ be a prompt and $y$ a response. At response position $t$ let $p=p(\cdot\mid x,y_{<t})$ be the teacher's and
$q=q_\theta(\cdot\mid x,y_{<t})$ the student's next-token distribution over the vocabulary $V$ ($|V|=49{,}152$).
Losses are averaged over response tokens. All equations are unit-tested against hand calculations on a 4-token
vocabulary.

\subsection{Teacher-signal formats (arms)}\label{sec:arms}
Let $T$ be the teacher's $k$ most probable tokens at a position, $p(T)=\sum_{i\in T}p_i$, and
$\hat p_i=p_i/p(T)$ for $i\in T$.
\begin{itemize}
\item \textbf{SFT}: $-\log q(y_t)$ on the gold response. \textbf{Seq-KD} \citep{kim2016seqkd}: the same on the
  teacher's greedy response.
\item \textbf{Full KD} (online): FKL $\sum_{i\in V}p_i\log\frac{p_i}{q_i}$; RKL $\sum_{i\in V}q_i\log\frac{q_i}{p_i}$.
\item \textbf{Top-$k$}, $k\in\{8,32,128\}$: FKL $\sum_{i\in T}\hat p_i\log\frac{\hat p_i}{q_i}$ with the student's
  \emph{full} softmax $q$, so the zero target outside $T$ pushes the student's tail mass towards zero. RKL
  renormalises both sides on $T$: $\KL(\hat q\Vert\hat p)=\sum_{i\in T}\hat q_i\log\frac{\hat q_i}{\hat p_i}$ with
  $\hat q_i=q_i/q(T)$; it is invariant to the student's logits outside $T$ and gives no gradient to its tail mass.
  Top-$k$ RKL is run at $k\in\{8,128\}$.
\item \textbf{Tail bucket} (ghost token; \citealp{anshumann2025sparse,dasgupta2026tad,huang2026taopd}), $k=8$:
  $(k{+}1)$-way distributions $\tilde p=(p_T,\,1-p(T))$ and $\tilde q=(q_T,\,1-q(T))$; FKL $\KL(\tilde p\Vert\tilde q)$,
  RKL $\KL(\tilde q\Vert\tilde p)$. It costs one extra number per position.
\item \textbf{RS-KD} \citep{anshumann2025sparse}: $N=50$ tokens drawn from the teacher at temperature 1,
  $\hat p_i=c_i/N$ with draw counts $c_i$; FKL $\KL(\hat p\Vert q)$ with the student's full softmax. A
  support-restricted reverse KL is biased (a once-drawn token receives target mass $\ge 1/N$), so RS-KD is run with
  forward KL only.
\item \textbf{Uniform tail} (allocation diagnostic, \S\ref{sec:tail}), $k=8$, FKL only: target $p_i$ on $T$ and the
  exact tail mass $1-p(T)$ spread uniformly over the other $|V|-k$ tokens.
\end{itemize}

\subsection{Cache and alignment}\label{sec:cache}
One teacher pass over the training responses stores, per response position, the top-128 probabilities (fp32),
their token ids, the exact leftover mass $1-p(T)$ computed in log space (no $1-\sum$ cancellation), and the 50
RS-KD draws. Smaller $k$ are slices of the same cache with the tail mass of ranks $k{+}1..128$ added back. The
teacher sees its chat template and the student the Alpaca/MiniLLM prompt; the \emph{same response token ids} are
appended to both (the models share one tokenizer), so the cache is indexed by exactly the response positions the
student is trained on. The cache is keyed by each example's prompt prefix and response ids and verified against
the training encodings before every run. \textbf{Online full KD uses the same prefixes, the same response masks
and the same positions}: the teacher runs on the identical teacher-side encoding during training, and its targets
go through the same top-$k$/slice code path as the cache, so an arm trained online and the same arm trained from
the cache differ only by floating-point rounding (checked in tests).

\subsection{Why a tail bucket restores amount but not identity}\label{sec:decomp}
For any set $T$, forward KL decomposes exactly (chain rule; the decomposition behind TAD, \citealp{dasgupta2026tad}):
\begin{equation}
\begin{aligned}
\KL(p\Vert q)={}&\KL(\tilde p\Vert\tilde q)\\
&+p(\bar T)\,\KL\big(p_{\mid\bar T}\,\Vert\,q_{\mid\bar T}\big),
\end{aligned}
\label{eq:decomp}
\end{equation}
where $\bar T=V{\setminus}T$ and $p_{\mid\bar T}=p(\cdot\mid\bar T)$ renormalises the teacher over the tail
(likewise $q_{\mid\bar T}$). The tail bucket optimises the first term
only: it matches \emph{how much} mass the student puts in the tail but leaves \emph{which} tail tokens receive it
untrained. Plain top-$k$ FKL optimises neither tail quantity; RS-KD's sampled targets include tail tokens
individually and estimate both terms. We therefore measure tail \emph{amount} (student mass outside the teacher's
top-8) and tail \emph{identity} separately. \textbf{Within-tail KL} is the last factor of Eq.~\ref{eq:decomp},
computed at every held-out position from the full teacher and student logits: both are renormalised over
$V{\setminus}T_8$ (masking the teacher's top-8 ids, log-sum-exp over the rest) and the KL is averaged over
positions. The amount/identity split was chosen after seeing seed 1; seeds 2--3 serve as its confirmation.

% =====================================================================================================
\section{Experimental setup}\label{sec:setup}

\paragraph{Models and data.} Teacher: SmolLM2-1.7B-Instruct (bf16). Student: SmolLM2-135M (base). Data:
databricks-dolly-15k following MiniLLM's protocol (prompts over 256 student tokens dropped, sequences cut at 512
tokens): 500 test and 1,000 validation examples, and 11,221 training examples after removing training items that
share an instruction or context with validation or test. Out-of-distribution test sets: Self-Instruct (242) and
Vicuna (80) prompts from MiniLLM's release.

\paragraph{Training and learning-rate selection.} Identical for every arm: AdamW (weight decay 0.01), batch 16,
3 epochs (2,106 steps), 5\% warm-up then linear decay, gradient clipping 1.0, fp32 weights with bf16 autocast and
an fp32 output layer, temperature 1, pure KD loss (no added cross-entropy). The learning rate was chosen once, on
validation only: SFT and full-KD FKL were trained on 2,000 examples at $\{1,2,3\}\times10^{-4}$ and the rate with
the best mean greedy validation ROUGE-L ($2\times10^{-4}$) was frozen for every arm before the main grid. Seeds 1,
2, 3. One RTX 4060 laptop GPU (8\,GB).

\paragraph{Evaluation.}\label{sec:eval}
\textbf{Reference overlap}: ROUGE-L (MiniLLM normalisation) of the greedy response and the mean over 8 samples.
\textbf{Diversity}: 8 samples per prompt at temperature 1 with no top-$k$/top-$p$ truncation (which would hide
exactly the tail we study): semantic diversity (mean pairwise cosine distance of all-MiniLM-L6-v2 embeddings),
Self-BLEU and distinct-2. \textbf{Teacher compatibility}: the teacher's mean log-likelihood per token of the
student's samples (reference-free). \textbf{Calibration} (teacher-forced on held-out gold responses): at each
position the student's confidence is its maximum next-token probability and it is correct if the argmax equals
the \emph{gold next token}; ECE is the position-weighted mean $|\text{accuracy}-\text{confidence}|$ over 15
equal-width confidence bins; we also report accuracy, NLL of the gold token and the multi-class Brier score
(Table~\ref{tab:calibration}, appendix). \textbf{Rare-token loss}: the student's NLL of gold tokens that occur fewer
than 10 times in the training responses and not in the prompt. \textbf{Long-tail knowledge} (exploratory): PopQA
accuracy by popularity quintile on the 2,444 questions the teacher answers correctly.

\paragraph{Endpoints and statistics.} Main contrast: top-8 vs.\ full KD, forward KL, Dolly test. We report mean
$\pm$ std over training seeds, per-seed effects, and a two-level (seed $\times$ prompt) paired bootstrap; ``no
difference'' in greedy ROUGE-L is claimed only within a $\pm1.0$ equivalence margin (90\% interval, two one-sided
tests).

% =====================================================================================================
\section{Results}

\subsection{Top-$k$ truncation reduces sampled diversity despite similar greedy ROUGE-L}\label{sec:main}

\input{\resultsdir tables/main_test}

Table~\ref{tab:main} gives the main grid. Relative to full KD, the top-8 student's greedy ROUGE-L differs by
\DiffRougeG{} points (90\% interval \GreedyCI{}, \GreedyEquiv{} the $\pm1.0$ margin), while its semantic diversity
is \TopEightSemDiv{} against \FullSemDiv{} (\SemDivDropPct\% lower), Self-BLEU is \TopEightSelfBleu{} against
\FullSelfBleu{}, ECE is \TopEightECE{} against \FullECE{} ($\times$\ECERatio), and rare-token loss is \TopEightRare{}
against \FullRare{}. Sampled quality is \emph{not} matched: top-8 samples have higher reference overlap
(\DiffRougeS{}) and higher teacher compatibility (\DiffCompat{} nats/token), i.e.\ the top-8 student is sharper;
\S\ref{sec:temp} therefore matches the samples' reference overlap before comparing diversity. Every metric moves
monotonically with $k$ (Figure~\ref{fig:dose}), and seed-to-seed standard deviations of semantic diversity (at most
\MaxSemDivStd{}) are an order of magnitude smaller than the top-8 vs.\ full gap (Table~\ref{tab:effects}). For
reference, the teacher's semantic diversity is \TeacherSemDiv{} and the untrained base student's \BaseSemDiv{}:
a 135M student is more dispersed than its teacher, so higher diversity is not automatically better.

\begin{figure}[t]\centering
\figorbox{dose}{\columnwidth}
\caption{Dose--response in the cached $k$ (forward KL, Dolly test, mean $\pm$ std over seeds), with full KD,
RS-KD and the tail bucket as reference lines.}\label{fig:dose}
\end{figure}

\input{\resultsdir tables/effects_test}

\subsection{The effect is not an artefact of answer length}\label{sec:length}
Top-8 samples are not systematically shorter in a way that explains the gap: on the first 32 words of every sample
\LenKeptFirst\% of the semantic-diversity gap remains, and on finished samples only \LenKeptFinished\%
(Table~\ref{tab:length}, appendix).

\subsection{Matched temperatures}\label{sec:temp}

\paragraph{Pre-registered design.} Before any of these runs, we fixed the following in the experiment log. On the
first 300 validation prompts (seed-1 students, 8 samples), top-8 is sampled at $T\in\{1.0,1.15,1.3,1.5\}$ and full
KD at $T\in\{0.7,0.85,1.0\}$. $T^\star_{\text{top-8}}$ is the top-8 temperature whose reference overlap is closest
to full KD at $T{=}1$, and $T^\star_{\text{full}}$ the full-KD temperature closest to top-8 at $T{=}1$; if the target
is not bracketed one extra step of 0.15 is added once, and nothing is extrapolated. On Dolly test (training seeds
1--2) we compare \textbf{A} (primary): top-8 at $T^\star_{\text{top-8}}$ vs.\ full KD at 1.0, and \textbf{B}: top-8 at
1.0 vs.\ full KD at $T^\star_{\text{full}}$. Per prompt we average each arm's 8-sample value over the two seeds and
take the paired difference; 10,000 bootstrap resamples of the 500 prompt IDs give 95\% intervals. Reference
overlap is equivalent if its interval lies within $\pm1.0$. Outcome (a), a \emph{residual deficit}: quality
equivalent and the semantic-diversity interval entirely below $-0.010$. Outcome (b), \emph{recovered by
temperature}: quality equivalent and the interval inside $\pm0.010$. Otherwise inconclusive. The headline is A's
verdict unless B gives the opposite conclusive verdict.

\input{\resultsdir tables/tempsel}

\paragraph{Result.} The selection gave $T^\star_{\text{top-8}}=\TTopStar$ and $T^\star_{\text{full}}=\TFullStar$
(Table~\ref{tab:tempsel}). In comparison A reference overlap is equivalent (difference \MatchAReferenceOverlap{},
95\% CI \MatchAReferenceOverlapCI{}) and the semantic-diversity difference is \MatchASemanticDiversity{}
\MatchASemanticDiversityCI{}; comparison B gives \MatchBReferenceOverlap{} \MatchBReferenceOverlapCI{} and
\MatchBSemanticDiversity{} \MatchBSemanticDiversityCI{}. Table~\ref{tab:matched} and Figure~\ref{fig:temp} give all
secondary measures (Self-BLEU, teacher compatibility, finish rate, length) and the sensitivity margins.

\paragraph{Interpretation.} Outcome (a), a residual diversity deficit, is \textbf{not supported}: in both comparisons
the diversity interval lies entirely \emph{above} zero, i.e.\ at matched reference overlap the top-8 student is more
diverse than full KD, and Self-BLEU agrees. Because the pre-registered rule defines only a deficit and equivalence
within $\pm0.010$, an interval above $+0.010$ is formally inconclusive. Part of the excess is explained by the small
remaining quality difference (top-8 has slightly lower reference overlap in both comparisons). In an
\emph{exploratory, post-hoc} adjustment (not pre-registered) we subtract full KD's own diversity-per-ROUGE-L slope on
test (\AdjSlope{}, chord between its $T{=}0.85$ and $T{=}1$ points) times the reference-overlap difference: the
adjusted diversity difference is \AdjA{} for A and \AdjB{} for B, i.e.\ approximately zero. At the level of
generations, top-$k$ caching therefore shifts the student along the quality--diversity frontier (it is sharper)
rather than below it; the teacher-forced differences in calibration and tail identity (\S\ref{sec:tail}) are not
affected by sampling temperature.

\begin{figure*}[t]\centering
\figorbox{temperature}{0.85\textwidth}
\caption{Left: validation temperature selection (seed 1), reference overlap vs.\ semantic diversity, labels are
sampling temperatures. Right: matched test comparisons A and B (top-8 minus full KD) with 95\% paired-bootstrap
intervals; shaded: the $\pm1.0$ quality and $\pm0.010$ diversity margins.}\label{fig:temp}
\end{figure*}

\input{\resultsdir tables/matched}

\subsection{The ghost token restores tail amount, not tail identity}\label{sec:tail}

\input{\resultsdir tables/recovery}

Table~\ref{tab:recovery} expresses each arm as the share of the top-8$\to$full gap it closes, per metric. The tail
bucket closes \RecTBAmount\% of the gap in tail amount (student mass outside the teacher's top-8) but only
\RecTBWithinTail\% of the gap in within-tail KL and \RecTBRare\% of the gap in rare-token loss; for semantic
diversity it closes \RecTBSemDiv\%, for Self-BLEU \RecTBSelfBleu\% and for ECE \RecTBECE\%. RS-KD closes
\RecRSAmount\% (tail amount), \RecRSWithinTail\% (within-tail KL), \RecRSRare\% (rare-token loss), \RecRSSemDiv\%
(semantic diversity), \RecRSSelfBleu\% (Self-BLEU) and \RecRSECE\% (ECE): its approximation fidelity extends to the
generation endpoints measured. This is the pattern Eq.~\ref{eq:decomp} predicts. The uniform-tail diagnostic (Table~\ref{tab:uniform}; \textbf{seed 1 only}) gives the student the exact tail
\emph{amount} but spread uniformly, i.e.\ with no tail identity. It raises tail amount (\UniAmount{} vs.\
\UTopAmount{} for top-8 and \UFullAmount{} for full KD), raises semantic diversity beyond full KD (\UniSemDiv{} vs.\
\UFullSemDiv{}) and lowers ECE (\UniECE{} vs.\ \UFullECE{}), but leaves within-tail KL where top-8 has it
(\UniWithinKL{} vs.\ \UTopWithinKL{}; full KD \UFullWithinKL{}) and only partly improves rare-token loss (\UniRare{}
vs.\ \UTopRare{}; full KD \UFullRare{}). Diversity and calibration thus respond to the amount of tail mass;
within-tail KL and, largely, rare-token prediction need the teacher's tail identity.

\subsection{Forward vs.\ reverse KL}\label{sec:rkl}
From the base checkpoint, reverse-KL students do not learn to end their responses: the finish rate of full-KD RKL
samples is \FinishRKLBase{} (Table~\ref{tab:main}), so their diversity numbers are not comparable and we report them
as a documented failure. We therefore test the divergence interaction from the same-seed SFT student
(``MiniLLM-style'' initialisation) in a matched $2\times4$ design (FKL/RKL $\times$ top-8, top-128, tail bucket,
full; Table~\ref{tab:ablation-sft}; seeds 1--2). There the full-KD RKL finish rate is \FinishRKLSFT{}, but top-8
RKL still does not terminate (finish rate \FinishTopEightRKLSFT{}). Its semantic diversity (\SFTTopEightRKLSemDiv{}
vs.\ \SFTFullRKLSemDiv{} for full-KD RKL) is therefore measured on unfinished, length-capped samples and is
\emph{not comparable}; we do not draw a top-$k$ $\times$ divergence conclusion from it. Under FKL from SFT
initialisation the main-grid pattern repeats (top-8 \SFTTopEightSemDiv{} vs.\ full KD \SFTFullSemDiv{}).

\subsection{Cross-entropy, out-of-distribution prompts and PopQA}
Adding gold-token cross-entropy (weight 0.5; Table~\ref{tab:ablation-ce}, seeds 1--2) makes both students much
more competent (greedy ROUGE-L \CETopEightRougeG{} for top-8 and \CEFullRougeG{} for full KD, against about 18.6 with
pure KD), and the gaps persist: semantic diversity \CETopEightSemDiv{} vs.\ \CEFullSemDiv{} and ECE \CETopEightECE{}
vs.\ \CEFullECE{}. On Self-Instruct and Vicuna (Table~\ref{tab:external}) the ordering of semantic diversity is the
same as on Dolly (top-8 $<$ top-128 $\approx$ tail bucket $<$ RS-KD $\approx$ full KD; e.g.\ top-8 \SITopEightSemDiv{}
vs.\ full \SIFullSemDiv{} on Self-Instruct, \VicTopEightSemDiv{} vs.\ \VicFullSemDiv{} on Vicuna). The PopQA probe
(Table~\ref{tab:popqa}) shows no difference between cache formats: Dolly carries few rare facts, so this null is
uninformative about long-tail \emph{knowledge}; the fact-independent rare-token loss above is our long-tail
measurement.

% =====================================================================================================
\section{Conclusion}
Caching only the teacher's top-$k$ probabilities is not neutral for what a distilled student generates. In one
model family on one instruction dataset, a top-8 cache made the student sharper: less diverse at the default
temperature, worse calibrated and worse on rare tokens, at equivalent greedy ROUGE-L. The sharpening is not a fixed
loss of diversity: at matched reference overlap a residual deficit was not supported. What the cache loses is the
identity of the tail. Restoring only the tail amount, with a ghost token or a uniform spread, recovers diversity and
calibration but not within-tail KL and only partly rare-token prediction; random-sampling distillation was the only
cached format that matched full KD across the endpoints we measured.

% =====================================================================================================
\section*{Limitations}
\begin{itemize}
\item \textbf{One model family, one training set.} SmolLM2-1.7B-Instruct $\to$ SmolLM2-135M, trained on Dolly only;
  all claims are scoped to this setting. We do not know whether the effects grow or shrink with student size or
  other families.
\item \textbf{No task with several verifiable answers.} Diversity is measured by embedding distance and $n$-gram
  overlap; we do not show that the extra diversity of full KD is \emph{useful}, and higher diversity is not
  automatically better for a 135M student that is more dispersed than its teacher.
\item \textbf{Reverse KL from the base checkpoint fails} (responses do not terminate); the divergence interaction is
  tested only from SFT initialisation, with two seeds.
\item \textbf{PopQA is null} and uninformative about long-tail knowledge on a dataset with few rare facts.
\item \textbf{Post-hoc split.} The tail amount/identity split was chosen after seeing seed 1 (seeds 2--3 confirm
  it), and the primary endpoints were declared after seeds 1--2.
\item \textbf{The matched comparison} uses temperatures selected on one seed, two training seeds on test, and a
  prompt-level bootstrap that conditions on those seeds (seed variance is not in its interval).
\item One learning rate, tuned on SFT and forward KL; reverse-KL arms may be disadvantaged.
\item \textbf{Calibration} differences may be largely fixable by post-hoc temperature scaling of the student's
  logits, which we did not test.
\item The quality adjustment of \S\ref{sec:temp} is exploratory and post-hoc.
\end{itemize}

\section*{Ethics and reproducibility}
All data and models are public. Code, configurations, every run (including failed runs) and the pre-registration
log are released with the paper.

\IfFileExists{acl.sty}{}{\bibliographystyle{plainnat}}% acl.sty sets acl_natbib itself
\bibliography{references}

% =====================================================================================================
\appendix
\section{Additional tables}\label{sec:appendix}
\input{\resultsdir tables/calibration}
\input{\resultsdir tables/length}
\input{\resultsdir tables/external}
\input{\resultsdir tables/ablation_ce}
\input{\resultsdir tables/ablation_sft}
\input{\resultsdir tables/diag_uniform}
\input{\resultsdir tables/popqa}

\begin{figure}[t]\centering
\figorbox{reliability_test}{\columnwidth}
\caption{Reliability diagram (Dolly test, seed 1, forward KL): confidence (maximum next-token probability) vs.\
accuracy (argmax = gold token), bins with at least 30 positions.}\label{fig:reliability}
\end{figure}

\section{Run accounting}
Every trained run, evaluation and failure is recorded in \texttt{runs/results.jsonl}, \texttt{runs/evals.jsonl} and
\texttt{runs/failed.jsonl}; the experiment log lists every incident and decision with its date. The temperature
selection and matched comparisons were pre-registered on 19 October 2026 before any of their runs (commit
57a44f4).

\end{document}
